\subsection{Figures}
\begin{figure}[!h]
	\centering
	\includesvg[width=.99\textwidth]{figs/base/calibrated/abm_growth_fig.svg}
 \caption{ Graphical representation of the forecasting performance of forecasting the growth rate and inflation from 2014:Q1-2016:Q4 of the ABM (red line), the VAR model (green line) compared with observed Eurostat data for Netherlands (black line). The 90 per cent confidence interval is depicted by the shaded area around the mean trajectory. Model results are obtained as an average of 100 Monte Carlo simulations.}
   \label{fig:gvamacro_growthinflation_yq}
\end{figure}

\begin{figure}[!h]
	\centering
	\includesvg[width=\textwidth]{figs/base/calibrated/abm_yearly_fig.svg}
	\caption{Graphical representation of the forecasting performance from 2014:Q1-2016:Q4 of the ABM (red line), the AR model (green line) compared with observed Eurostat data for Netherlands (black line). The variables GDP, household consumption, investmentment, government consumption, and imports are measured in real terms. The 90 per cent confidence interval is depicted by the shaded area around the mean trajectory. Model results are obtained as an average of 100 Monte Carlo simulations.}     \label{fig:gvamacro_gvaetal_y}
\end{figure}

\begin{figure}[!h]
	\centering
	\includesvg[width=\textwidth]{figs/base/calibrated/abm_quarterly_fig.svg}
 \caption{Graphical representation of the forecasting performance from 2014:Q1-2016:Q4 of the ABM (red line), the AR model (green line) compared with observed Eurostat data for Netherlands (black line). The variables GDP, household consumption, investmentment, government consumption, and imports are measured in real terms. The 90 per cent confidence interval is depicted by the shaded area around the mean trajectory. Model results are obtained as an average of 100 Monte Carlo simulations.}
  \label{fig:gvamacro_gvaetal_y}
\end{figure}

\begin{figure}[!h]
	\centering
	\includesvg[width=.99\textwidth]{figs/base/calibrated/abmx_growth_fig.svg}
 \caption{ Graphical representation of the forecasting performance of forecasting the growth rate and inflation from 2014:Q1-2016:Q4 of the conditional ABM (red line), the VARX model (green line) compared with observed Eurostat data for Netherlands (black line). The 90 per cent confidence interval is depicted by the shaded area around the mean trajectory. Model results are obtained as an average of 100 Monte Carlo simulations.}
   \label{fig:gvamacro_growthinflation_yq}
\end{figure}

\begin{figure}[!h]
	\centering
	\includesvg[width=\textwidth]{figs/base/calibrated/abmx_yearly_fig.svg}
	\caption{Graphical representation of the forecasting performance from 2014:Q1-2016:Q4 of the conditional ABM (red line), the ARX model (green line) compared with observed Eurostat data for Netherlands (black line). The variables GDP, household consumption, investmentment, government consumption, and imports are measured in real terms. The 90 per cent confidence interval is depicted by the shaded area around the mean trajectory. Model results are obtained as an average of 100 Monte Carlo simulations.}     \label{fig:gvamacro_gvaetal_y}
\end{figure}

\begin{figure}[!h]
	\centering
	\includesvg[width=\textwidth]{figs/base/calibrated/abmx_quarterly_fig.svg}
 \caption{Graphical representation of the forecasting performance from 2014:Q1-2016:Q4 of the conditional ABM (red line), the ARX model (green line) compared with observed Eurostat data for Netherlands (black line). The variables GDP, household consumption, investmentment, government consumption, and imports are measured in real terms. The 90 per cent confidence interval is depicted by the shaded area around the mean trajectory. Model results are obtained as an average of 100 Monte Carlo simulations.}
  \label{fig:gvamacro_gvaetal_y}
\end{figure}


\begin{figure}[!h]
	\centering
	\includesvg[width=.99\textwidth]{figs/extended_heuristic/calibrated/abm_growth_fig.svg}
 \caption{ Graphical representation of the forecasting performance of forecasting the growth rate and inflation from 2014:Q1-2016:Q4 of the EH-ABM (red line), the VAR model (green line) compared with observed Eurostat data for Netherlands (black line). The 90 per cent confidence interval is depicted by the shaded area around the mean trajectory. Model results are obtained as an average of 100 Monte Carlo simulations.}
   \label{fig:gvamacro_growthinflation_yq}
\end{figure}

\begin{figure}[!h]
	\centering
	\includesvg[width=\textwidth]{figs/extended_heuristic/calibrated/abm_yearly_fig.svg}
	\caption{Graphical representation of the forecasting performance from 2014:Q1-2016:Q4 of the EH-ABM (red line), the AR model (green line) compared with observed Eurostat data for Netherlands (black line). The variables GDP, household consumption, investmentment, government consumption, and imports are measured in real terms. The 90 per cent confidence interval is depicted by the shaded area around the mean trajectory. Model results are obtained as an average of 100 Monte Carlo simulations.}     \label{fig:gvamacro_gvaetal_y}
\end{figure}

\begin{figure}[!h]
	\centering
	\includesvg[width=\textwidth]{figs/extended_heuristic/calibrated/abm_quarterly_fig.svg}
 \caption{Graphical representation of the forecasting performance from 2014:Q1-2016:Q4 of the EH-ABM (red line), the AR model (green line) compared with observed Eurostat data for Netherlands (black line). The variables GDP, household consumption, investmentment, government consumption, and imports are measured in real terms. The 90 per cent confidence interval is depicted by the shaded area around the mean trajectory. Model results are obtained as an average of 100 Monte Carlo simulations.}
  \label{fig:gvamacro_gvaetal_y}
\end{figure}

\begin{figure}[!h]
	\centering
	\includesvg[width=.99\textwidth]{figs/extended_heuristic/calibrated/abmx_growth_fig.svg}
 \caption{ Graphical representation of the forecasting performance of forecasting the growth rate and inflation from 2014:Q1-2016:Q4 of the conditional EH-ABM (red line), the VARX model (green line) compared with observed Eurostat data for Netherlands (black line). The 90 per cent confidence interval is depicted by the shaded area around the mean trajectory. Model results are obtained as an average of 100 Monte Carlo simulations.}
   \label{fig:gvamacro_growthinflation_yq}
\end{figure}

\begin{figure}[!h]
	\centering
	\includesvg[width=\textwidth]{figs/extended_heuristic/calibrated/abmx_yearly_fig.svg}
	\caption{Graphical representation of the forecasting performance from 2014:Q1-2016:Q4 of the conditional EH-ABM (red line), the ARX model (green line) compared with observed Eurostat data for Netherlands (black line). The variables GDP, household consumption, investmentment, government consumption, and imports are measured in real terms. The 90 per cent confidence interval is depicted by the shaded area around the mean trajectory. Model results are obtained as an average of 100 Monte Carlo simulations.}     \label{fig:gvamacro_gvaetal_y}
\end{figure}

\begin{figure}[!h]
	\centering
	\includesvg[width=\textwidth]{figs/extended_heuristic/calibrated/abmx_quarterly_fig.svg}
 \caption{Graphical representation of the forecasting performance from 2014:Q1-2016:Q4 of the conditional EH-ABM (red line), the ARX model (green line) compared with observed Eurostat data for Netherlands (black line). The variables GDP, household consumption, investmentment, government consumption, and imports are measured in real terms. The 90 per cent confidence interval is depicted by the shaded area around the mean trajectory. Model results are obtained as an average of 100 Monte Carlo simulations.}
  \label{fig:gvamacro_gvaetal_y}
\end{figure}
